\SourceLeaf{39}

Iam quoniam quaeque permutatio ex $p$ rebus constat, patet cuivis $p-1$ similes adinveniri posse, si ea res quae prima fuerat, ad secundum, tertium etc. locum promoveatur. Quarum si nullae identicae esse possunt manifestum est, omnium permutationum numerum per $p$ divisibilem evadere, quippe qui $p$ vicibus maior sit quam numerus omnium permutationum dissimilium. Supponamus igitur duas permutationes

\[PQ\ldots TV\ldots YZ;\qquad V\ldots YZPQ\ldots T\]

\noindent quarum altera ex altera per terminorum promotionem orta sit, identicas esse sive $P=V$ etc. Sit terminus $P$ qui in priori est primus, $n+1^{\mathrm{tus}}$ in posteriori. Erit igitur in serie posteriori terminus $n+1^{\mathrm{tus}}$ aequalis primo, $n+2^{\mathrm{tus}}$ secundo etc. unde $2n+1^{\mathrm{tus}}$ rursus primo aequalis evadet, eademque ratione $3n+1^{\mathrm{tus}}$ etc.; generaliterque terminus $kn+m^{\mathrm{tus}}$ $m^{\mathrm{to}}$ (ubi quando $kn+m$ ipsum $p$ superat, aut series $V\ldots YZPQ\ldots T$ semper ab initio repeti concipienda est, aut a $kn+m$ multiplum ipsius $p$ proxime minus rescindendum). Quamobrem si $k$ ita determinatur, ut fiat $kn\equiv1\ (\mathrm{mod.}p)$, quod fieri potest quia $p$ primus, sequitur generaliter terminum $m^{\mathrm{tum}}$ $m+1^{\mathrm{to}}$ aequalem esse, sive quemvis terminum sequenti, \emph{i. e.} omnes terminos aequales esse contra hypothesin.

\Article{42}

\emph{Si coefficientes $A,B,C\ldots.N$; $a,b,c\ldots.n$ duarum functionum formae}

\[
\begin{array}{l@{\qquad}c}
x^m+Ax^{m-1}+Bx^{m-2}+Cx^{m-3}\mathinner{\ldotp\ldotp\ldotp\ldotp\ldotp}+N
\quad\mathinner{\ldotp\quad\ldotp\quad\ldotp\quad\ldotp\quad\ldotp}&(P)\\
x^\mu+ax^{\mu-1}+bx^{\mu-2}+cx^{\mu-3}\mathinner{\ldotp\ldotp\ldotp\ldotp\ldotp}+n
\quad\mathinner{\ldotp\quad\ldotp\quad\ldotp\quad\ldotp\quad\ldotp}&(Q)
\end{array}
\]

\noindent\emph{omnes sunt rationales, neque vero omnes integri, productumque ex $(P)$ et $(Q)$}

\[=x^{m+\mu}+\mathfrak{A}x^{m+\mu-1}+\mathfrak{B}x^{m+\mu-2}+\emph{etc.}+\mathfrak{Z}\]

\noindent\emph{omnes coefficientes $\mathfrak{A},\mathfrak{B}\ldots.\mathfrak{Z}$ integri esse nequeunt.}

\emph{Demonstr.} Exprimantur omnes fractiones in coefficientibus $A,B$ etc. $a,b$ etc. per numeros quam minimos, eligaturque ad libitum numerus primus $p$, qui aliquem aut plures ex denominatoribus harum fractionum metiatur. Ponamus, id quod licet, $p$ metiri denominatorem alicuius coefficientis fracti in $(P)$, patetque si $(Q)$ per $p$ dividatur, etiam in $\frac{(Q)}{p}$ dari ad minimum unum coefficientem fractum cuius denominator implicet factorem $p$ (puta coefficientem primum $\frac{1}{p}$).

\SourceLeaf{40}

Iam facile perspicitur, in $(P)$ datum iri terminum unum, fractum, cuius denominator involvat \emph{plures} dimensiones ipsius $p$ quam denominatores omnium similium praecedentium, et \emph{non pauciores} quam denominatores omnium sequentium; sit hic terminus $=Gx^g$, et multitudo dimensionum ipsius $p$ in denominatore ipsius $G,=t$. Similis terminus dabitur in $\frac{(Q)}{p}$ qui sit $=\varGamma x^\gamma$ et multitudo dimensionum ipsius $p$ in denominatore ipsius $\varGamma,=\tau$. Manifesto hic erit $t+\tau$ ad minimum $=2$. His ita praeparatis, terminus $x^{g+\gamma}$ producti ex $(P)$ et $(Q)$ coefficientem habebit fractum, cuius denominator $t+\tau-1$ dimensiones ipsius $p$ involvet, id quod ita demonstratur.

Sint termini qui in $(P)$ terminum $Gx^g$ praecedunt, ${}'Gx^{g+1}$, ${}''Gx^{g+2}$ etc. sequentes vero $G'x^{g-1}$, $G''x^{g-2}$ etc.; similiterque in $\frac{(Q)}{p}$ praecedant terminum $\varGamma x^\gamma$ termini ${}'\varGamma x^{\gamma+1}$, ${}''\varGamma x^{\gamma+2}$ etc. sequantur autem termini $\varGamma'x^{\gamma-1}$, $\varGamma''x^{\gamma-2}$ etc. Tum constat in producto ex $(P)$, $\frac{(Q)}{p}$ coefficientem termini $x^{g+\gamma}$ fore

\[
\begin{aligned}
&=G\varGamma+{}'G\varGamma'+{}''G\varGamma''+\text{etc.}\\
&\phantom{=G\varGamma}\mathord{+}{}'\varGamma G'+{}''\varGamma G''+\text{etc.}
\end{aligned}
\]

Pars $G\varGamma$ erit fractio quae si per numeros quam minimos exprimitur in denominatore $t+\tau$ dimensiones ipsius $p$ involvit, reliquae autem partes si sunt fractae, in denominatore pauciores dimensiones numeri $p$ implicabunt, quoniam omnes sunt producta e binis factoribus quorum alter non plures quam $t$, alter vero pauciores quam $\tau$ dimensiones ipsius $p$ implicat; vel alter non plures quam $\tau$, alterque pauciores quam $t$. Hinc $G\varGamma$ erit formae $\frac{e}{fp^{t+\tau}}$ reliquarum vero summa formae $\frac{e'}{f'p^{t+\tau-\delta}}$ ubi $\delta$ positivus et $e,f,f'$ a factore $p$ liberi: quare omnium summa erit $=\frac{ef'+e'fp^\delta}{ff'p^{t+\tau}}$ cuius numerator per $p$ non divisibilis, adeoque denominator per nullam reductionem pauciores dimensiones quam $t+\tau$ obtinere potest. Hinc coefficiens termini $x^{g+\gamma}$ in producto ex $(P),(Q)$ erit

\[=\frac{ef'+e'fp^\delta}{ff'p^{t+\tau-1}}.\]

\noindent\emph{i. e. fractio} cuius denominator $t+\tau-1$ dimensiones ipsius $p$ implicat.\par
\noindent\mbox{\emph{Q. E. D.}}

\SourceLeaf{41}

\Article{43}

\emph{Congruentia $m^{ti}$ gradus}

\[Ax^m+Bx^{m-1}+Cx^{m-2}+\emph{etc.}+Mx+N\equiv0\]

\noindent\emph{cuius modulus est numerus primus $p$, ipsum $A$ non metiens, pluribus quam $m$ modis diversis solvi non potest, sive plures quam $m$ radices secundum $p$ incongruas non habet} (Vid. artt. 25, 26).

Si quis neget, ponamus dari congruentias diversorum graduum $m,n$ etc. quae plures quam $m,n$ etc. radices habeant, sitque minimus gradus $m$, ita ut omnes similes congruentiae inferiorum graduum theoremati nostro sint consentaneae. Quod quum de primo gradu iam supra sit demonstratum (art. 26), manifestum est, $m$ fore aut $=2$ aut maiorem. Admittet itaque congruentia

\[Ax^m+Bx^{m-1}+\text{etc.}+Mx+N\equiv0\]

\noindent saltem $m+1$ radices, quae sint $x\equiv\alpha,x\equiv\beta,x\equiv\gamma$ etc., ponamusque id quod licet omnes numeros $\alpha,\beta,\gamma$ etc. esse positivos et minores quam $p$, omniumque minimum $\alpha$. Iam in congruentia proposita substituatur pro $x$, $y+\alpha$, transeatque inde in hanc

\[A'y^m+B'y^{m-1}+C'y^{m-2}....+M'y+N'\equiv0\]

Tum manifestum est, huic congruentiae satisfieri, si ponatur $y\equiv0$, aut $\equiv\beta-\alpha$, aut $\equiv\gamma-\alpha$ etc., quae radices omnes erunt diversae, numerusque earum $=m+1$. At ex eo quod $y\equiv0$ est radix, sequitur, $N'$ per $p$ divisibilem fore. Quare etiam haec expressio

\[y(A'y^{m-1}+B'y^{m-2}+\text{etc.}+M')\text{ fiet }\equiv0\ (\mathrm{mod.}p),\]

\noindent si ipsi $y$ unus ex $m$ valoribus $\beta-\alpha,\gamma-\alpha$ etc. tribuitur, qui omnes sunt $>0$ et $<p$, adeoque in omnibus hisce casibus etiam

\[A'y^{m-1}+B'y^{m-2}+\text{etc.}+M'\text{ fiet }\equiv0\ \text{(art. 22)}\]

\noindent\emph{i. e.} congruentia \qquad $A'y^{m-1}+B'y^{m-2}+\text{etc.}+M'\equiv0$

\noindent quae est gradus $m-1^{\mathrm{ti}}$, $m$ radices habet et proin theoremati nostro adversatur (patet enim facile, $A'$ fore $=A$, adeoque per $p$ non divisibilem, uti requiritur) licet supposuerimus, omnes congruentias inferioris gradus quam $m^{\mathrm{ti}}$, theoremati consentire. \mbox{\emph{Q. E. A.}}

\SourceLeaf{42}

\Article{44}

Quamvis hic supposuerimus, modulum $p$ non metiri coefficientem termini summi, tamen theorema ad hunc casum non restringitur. Si enim primus coefficiens sive etiam aliqui sequentium per $p$ divisibiles essent, hi termini tuto reiici possent, congruentiaque tandem ad inferiorem gradum deprimeretur, ubi coefficiens primus per $p$ non amplius foret divisibilis, siquidem non omnes coefficientes per $p$ dividi possunt; in quo casu congruentia foret identica atque incognita prorsus indeterminata.

Theorema hoc primum ab ill. La Grange propositum atque demonstratum est (\emph{M\'em. de l'Ac. de Berlin, Ann\'ee} 1768 \emph{p.} 192). Exstat etiam in dissert. ill. Le Gendre, \emph{Recherches d'Analyse ind\'etermin\'ee, Hist. de l'Acad. de Paris} 1785 \emph{p.} 466. Ill. Euler in \emph{Nov. Comm. Ac. Petr.} XVIII \emph{p.} 93 demonstravit congruentiam $x^n-1\equiv0$ plures quam $n$ radices diversas habere non posse. Quae quamvis sit particularis, tamen methodus qua vir summus usus est omnibus congruentiis facile adaptari potest. Casum adhuc magis limitatum iam antea absolverat, \emph{Comm. nov. Ac. Petr.} V \emph{p.} 6, sed haec methodus generaliter adhiberi nequit. Infra Sect. VIII alio adhuc modo theorema demonstrabimus; at quantumvis diversae primo aspectu omnes hae methodi videri possint, periti qui comparare eas voluerint facile certiores fient omnes eidem principio superstructas esse. Ceterum quum hoc theorema hic tantum tamquam lemma sit considerandum, neque completa expositio huc pertineat: de modulis compositis seorsim agere supersedemus.

\BookRule{0.28\linewidth}
